% ----------------------------------------------------------------
\subsection{Le dimensioni del grano}\label{sec:dimensioni}

Sedimentati i parametri
precedenti, procederò sovrapponendone di nuovi, e dalla loro
combinazione nasce la prima texture complessa.

La posizione di un grano nel campo stereofonico è
codificata in Mid-Side: l'angolo \texttt{pan} bilancia una componente
centrale e una laterale, da cui si ottengono per somma e differenza i due canali,
\begin{equation}
  M = s\cos\theta, \quad S = s\sin\theta,
  \qquad L = \frac{M+S}{\sqrt{2}}, \quad R = \frac{M-S}{\sqrt{2}},
  \label{eq:midside}
\end{equation}
con $s$ il segnale del grano e $\theta$ l'angolo. La potenza si conserva
($L^2 + R^2 = s^2$); a $\theta = 0$ il grano è centrato, a $\pm 45^\circ$
è tutto su un canale, oltre i $\pm 45^\circ$ la componente laterale supera
la centrale e l'immagine eccede gli altoparlanti. È la matrice
somma-differenza di Blumlein~\cite{Blumlein1931}, dove la larghezza si
regola sull'ampiezza del lato. Come ogni parametro, \texttt{pan} ammette
una traiettoria centrale e un range: \texttt{pan\_range} distribuisce i
grani sul campo, grano per grano, secondo lo schema di
\S\ref{sec:deviazione}.

    \lstinputlisting[
      linerange=23,
      language=yaml,
      basicstyle=\ttfamily\fontsize{7.5}{9}\selectfont,
      label=lst:dimensioni,
      caption={Le dimensioni restanti all'opera insieme, sul puntatore
      tenuto quasi fermo (\texttt{ex5\_texture.yml}).}
    ]{examples/duration/duration.yml}

L'esempio le mette all'opera tutte insieme (Listato~\ref{lst:dimensioni}).
Sul puntatore tenuto quasi fermo, le dimensioni sedimentate e le nuove si
muovono di concerto: la densità oscilla fra estremi, la
\texttt{distribution} dissolve e riforma la griglia, la durata del grano
si contrae e si riapre, il \texttt{dephase} regola quanti grani
decorrelano, mentre \texttt{pan} e \texttt{pan\_range} aprono il campo. La
MAP (Fig.~\ref{fig:dimensioni}) restituisce questa sovrapposizione come
un'unica texture in evoluzione: non più un parametro isolato su uno sfondo
fermo, ma molte dimensioni che compongono insieme. Mancano ancora la
finestra nella sua forma piena e l'altezza: l'esempio completo le aggiunge
e spinge ogni envelope oltre.

\begin{figure}[h]
    \includegraphics[width=\linewidth]{examples/duration/duration_map.pdf}\\
    \captionof{figure}{La MAP del rendering (assi come in
    Fig.~\ref{fig:pointer-MAP}). Sul puntatore quasi fermo, densità,
    \texttt{distribution}, durata del grano, \texttt{dephase}, \texttt{pan}
    e \texttt{pan\_range} si muovono insieme: la sovrapposizione dei
    parametri sedimentati e dei nuovi compone un'unica texture in
    evoluzione.}
    \label{fig:dimensioni}
\end{figure}